\section{Galois-stable induction systems and canonical normal isotypes}
\label{sec:SF-galois-stable-normal-isotypes}

\paragraph{Author freeze and provenance.}
D-TD60, question Q2fp, is a separate extension of D-TD56,
not a change to that accepted normal-induction theorem. The parent requested a verifiable nonnormal
subclass; the author identified the canonical normal-subgroup isotypic
system and the exact normalizer criterion below. The complete proof
has passed the author's complete reading and complete root and
independent whole-proof reviews. No numerical execution or source
existence is asserted.
The original-field and natural-polarization inputs retain the exact
accepted D-OG-01--03, HD-PD01 and D-TD56 scopes.

\subsection*{Actual source and additional block hypothesis}

Use the actual source of D-TD56: a connected finite
\(G\)-Galois cover \(C\to\mathbf P^1\) in a family, with a faithful
absolutely irreducible representation \(\rho_K\) over its exact
imaginary quadratic character field \(K\). Its multiplicity variation
has \(K\)-dimension six, signature \((3,3)\), and actual period
rank nine. Fix the signed determinant class \(\delta\) of the
naturally polarized primitive sixfold selected by a rational
Rosati-fixed rank-one projector. The comparison class in any
minimality assertion includes nonprimitive sources with this exact
field, signature, period rank and natural determinant class.

Suppose that, as a complex representation,
\[
 \rho_{\mathbf C}=\operatorname{Ind}_H^G\tau,
 \qquad e=\dim\tau\ge2,
 \qquad k=[G:H],\qquad L=\mathbf Q(\chi_\tau),
 \qquad \ell=[L:K].
 \tag{GSB.1}
\]
The inducing representation \(\tau\) is irreducible, as is forced by
irreducibility of its induction. The subgroup \(H\) is not initially
assumed normal. Choose the corresponding direct-sum imprimitivity system
\[
 \mathcal B=\{gB:gH\in G/H\},
 \qquad \dim B=e,\qquad\operatorname{Stab}_G(B)=H.
 \tag{GSB.2}
\]
Here \(B\) carries the distinguished inducing representation
\(\tau\). The \(K\)-model gives coefficientwise action of
\(\operatorname{Gal}(\overline{\mathbf Q}/K)\) on algebraic
subspaces of \(\rho\).

\begin{theorem}[D-TD60: Galois-stable block reduction]
The following two conditions are equivalent:
\[
 \begin{split}
 &\operatorname{Gal}(\overline{\mathbf Q}/K)
       \text{ preserves the chosen set }\mathcal B;\\
 &\text{for every }\sigma\in\operatorname{Gal}(L/K),\quad
       \tau^\sigma\simeq\tau^{n_\sigma}
       \text{ for some }n_\sigma\in N_G(H).
 \end{split}
 \tag{GSB.3}
\]
Assume them. There is a subgroup \(H'\subseteq N_G(H)\), containing
\(H\), and an absolutely irreducible exact-\(K\) model \(W_K\)
such that
\[
 \begin{gathered}
 H\triangleleft H',\qquad [H':H]=\ell,\qquad
 H'/H\simeq\operatorname{Gal}(L/K),\\
 W_{\mathbf C}=\operatorname{Ind}_H^{H'}\tau,
 \qquad \operatorname{Ind}_{H'}^G W_K=\rho_K,
 \qquad \dim_KW_K=e\ell.
 \end{gathered}
 \tag{GSB.4}
\]
Put \(N=\ker\tau\). Then \(\ker W_K=N\triangleleft H'\), and the
actual connected cover
\[
 C/N\longrightarrow C/H',\qquad J'=H'/N,
 \qquad |J'|=\ell|H/N|
 \tag{GSB.5}
\]
is an exact-\(K\), signature-\((3,3)\), rank-nine source with natural
signed determinant \(\delta\). Both \(C/H\) and \(C/H'\) are rational.
If \(\ell>1\), the group \(H'/H\) is cyclic or Klein four, with
\[
 \begin{cases}
 \ell\le\lfloor(e+3)/5\rfloor,&\text{in the cyclic case},\\
 \ell=4,\quad e\ge15,&\text{in the Klein-four case}.
 \end{cases}
 \tag{GSB.6}
\]
If the original coefficient degree is minimal in the stated actual
source comparison class, then
\[
 H'=G,\qquad H\triangleleft G,\qquad k=\ell,
 \qquad N=1.
 \tag{GSB.7}
\]
Thus a genuinely nonnormal inducing system satisfying (GSB.3)
cannot occur in such a minimal source.
\end{theorem}

\subsection*{1. Algebraic blocks and the normalizer criterion}

\begin{proof}
Frobenius reciprocity and irreducibility give
\[
 \dim\operatorname{Hom}_H(\tau,\rho|_H)
   =\dim\operatorname{Hom}_G(\operatorname{Ind}_H^G\tau,\rho)=1.
 \tag{GSB.8}
\]
Consequently the central \(\tau\)-idempotent in
\(\overline{\mathbf Q}[H]\) has image exactly \(B\), not an
unspecified larger isotypic component. Character values are algebraic,
so \(B\) is algebraic. Induced-character values are sums of
\(\chi_\tau\)-values; hence \(K\subseteq L\). The character field
\(L\) is a subfield of a cyclotomic field, so \(L/\mathbf Q\),
and therefore \(L/K\), is Galois abelian.

For each coefficient automorphism fixing \(K\), induction commutes
with that automorphism and fixes \(\rho\). Thus the analogue of
(GSB.8) holds for every \(\tau^\sigma\). Its unique isotypic image
is \(\sigma(B)\). If the first condition of (GSB.3) holds, write
\(\sigma(B)=gB\). Since the representation matrices are over \(K\),
this space is stable under \(H\). Its full stabilizer is
\(gHg^{-1}\), so \(H\subseteq gHg^{-1}\). Equality follows from
equal finite orders; hence \(g\in N_G(H)\). Its \(H\)-character
is both \(\chi_{\tau^\sigma}\) and \(\chi_{\tau^g}\), proving
the second condition.

Conversely suppose the second condition. The space \(n_\sigma B\)
is an \(H\)-submodule of type \(\tau^\sigma\), and (GSB.8) for that
type makes it the unique such isotypic subspace. Therefore it is
\(\sigma(B)\). Galois commutes with every original \(G\)-matrix,
so it also preserves the translates of \(B\), proving the first
condition. This establishes exactly the stated equivalence, not an
automatic assertion for arbitrary imprimitivity systems.

An automorphism fixing \(K\) stabilizes \(B\) if and only if it fixes
\(\chi_\tau\): one direction takes traces on the invariant space,
and the other uses its central idempotent over \(L\) and (GSB.8).
Thus the orbit of \(B\) has size \(\ell\), and the action on it is
the regular action of \(\Gamma=\operatorname{Gal}(L/K)\). The action
on all \(\mathcal B\) is determined by this action, since it commutes
with the transitive \(G\)-action.

\subsection*{2. The centralizer of a transitive block action}

A permutation of \(G/H\) commuting with \(G\) is determined by its
value \(gH\) at \(H\). That value must be \(H\)-fixed, which means
\(H\subseteq gHg^{-1}\), hence \(g\in N_G(H)\). Conversely such
a coset gives the commuting opposite right translation. This proves
\[
 C_{\operatorname{Sym}(G/H)}(G)
          \simeq (N_G(H)/H)^{\mathrm{op}}.
 \tag{GSB.9}
\]
The centralizer acts semiregularly: a commuting permutation fixing one
block fixes all blocks. The Galois image is consequently a subgroup
of this centralizer isomorphic to \(\Gamma\), not merely a quotient
with an unspecified kernel. Let \(H'\) be its preimage in \(N_G(H)\).
Then \(H\triangleleft H'\), the quotient is the abelian group
\(\Gamma\), and the Galois orbit \(\mathcal O\) of \(B\) is exactly
\(\{h'B:h'\in H'\}\). Its stabilizer in \(G\) is \(H'\):
if \(g\mathcal O=\mathcal O\), then \(gB=h'B\) for some \(h'\),
so \(g\in h'H=H'\). Its \(G\)-translates partition \(\mathcal B\).

The sum of its central \(H\)-projectors is explicitly
\[
 p_{\mathcal O}=\frac{e}{|H|}\sum_{h\in H}
       \operatorname{Tr}_{L/K}(\chi_\tau(h^{-1}))\,h
       \in K[H],\qquad W_K=p_{\mathcal O}\rho_K.
 \tag{GSB.10}
\]
The distinct constituents \(\tau^\sigma\) each have multiplicity one.
Their \(H'\)-action is transitive, so \(W_K\) is absolutely
irreducible. The disjoint translated orbit sums, all defined over
\(K\), identify their direct sum with \(\rho_K\) and give the
induction identity in (GSB.4) over \(K\). Its character field is
contained in \(K\), and the induction identity conversely puts
\(K=\mathbf Q(\chi_\rho)\) inside that field. The model is thus
over its exact character field, with no Schur-index substitution.

All Galois-conjugate \(\tau\)-characters have the same kernel.
For a finite-order matrix of degree \(e\), its trace equals \(e\)
if and only if it is the identity; this follows by unitarity and
equality in the triangle inequality. This condition is unchanged
by a coefficient automorphism. Inside \(H\), the common kernel
of the orbit sum is therefore \(N\). An element of \(H'\setminus H\)
permutes the blocks of \(\mathcal O\) nontrivially in its regular
quotient action, and cannot act identically. This proves
\(\ker W_K=N\), including normality in \(H'\).

\subsection*{3. The actual family and its natural determinant}

The subgroup \(N\) need not be normal in \(G\); normality in
\(H'\) is exactly what is needed for (GSB.5). These are actual
intermediate quotient curves in the original family. Frobenius
reciprocity for \(\rho_K=\operatorname{Ind}_{H'}^G W_K\) identifies
the \(K\)-multiplicity local systems, not merely their dimensions
on one fibre. The quotient pullback identifies this system with
the \(W_K\)-multiplicity in \(H^1(C/N)\). Its signature and the
actual period rank on the original family remain \((3,3)\) and nine.
The period map factors through the actual quotient Hurwitz component,
whose period rank is therefore at least nine and hence exactly nine.
The natural rational determinant class is locally constant on that
component. Thus this is a replacement source in the stated family
class, not only a comparison at a single fibre.

For the natural polarization one uses exactly the rank-one
Rosati-projector comparison in D-TD56, Section 3, and the accepted
D-OG-01--03/HD-PD01 interface. The rational block for \(W_K\) is
\(\operatorname{Mat}_{e\ell}(K)\). A positive Rosati-fixed primitive
line, followed by the quotient-average projector for \(N\), has
support only on the original \(\rho\)-block, by the multiplicity-one
induction identity. It cuts out a rank-one coefficient line there.
Quotient pullback multiplies the Jacobian pairing by \(|N|\);
the two coefficient-line positive Hermitian norms have rational
ratio. Consequently the rank-six multiplicity Hermitian forms
differ by a positive rational scalar \(c\), and
\[
                  c^6=N_{K/\mathbf Q}(c^3).
 \tag{GSB.11}
\]
Their signed determinant classes coincide. This is a comparison of
the natural rational determinant classes, not equality of integral
polarization types or an inference from a norm in a larger field.

The accepted Euler/period calculation, repeated in D-TD56, Section 3,
applies to the complex coefficient \(\tau\) as well as to \(W_K\).
Each has complex multiplicity six and actual period rank nine, and
each degree is at least two. Genus one has at most six moving-rank
units and period rank at most six; genus two has cover-family
dimension at most five; larger base genus violates the nonnegative
moving-rank formula. Thus both bases \(C/H\) and \(C/H'\) are rational.
This step does not require a \(K\)-model for \(\tau\).

Now apply the already proved normal-induction theorem D-TD56 to
the actual faithful \(J'=H'/N\)-source, induced from its normal
subgroup \(H/N\) with faithful degree-\(e\) coefficient \(\tau\).
Its block index is \(\ell\) and its coefficient field is exactly
\(K\). That theorem gives the cyclic/Klein-four alternative and
both inequalities in (GSB.6), with its exact genuine-mark and
equal-kernel argument. There is no inference that the original
\(H\)-conjugate blocks outside \(H'\) have equal kernels.

Finally, if \(H'\ne G\), then \(e\ell<ek=\dim_K\rho_K\).
The actual replacement just constructed has all the prescribed
field, period, signature and natural-determinant data, contradicting
minimality. Hence \(H'=G\). Since \(H'\subseteq N_G(H)\), this
forces \(H\triangleleft G\) and \(k=\ell\). The equality
\(W_K=\rho_K\) makes \(N=1\) by original faithfulness. This proves
every assertion of the first theorem.
\end{proof}

\subsection*{4. The automatic canonical normal-subgroup case}

\begin{corollary}[D-TD60: canonical normal-isotypic reduction]
Let \(D\triangleleft G\). Choose a canonical isotypic component
\(B\) of \(\rho_{\mathbf C}|_D\), with irreducible \(D\)-character
\(\theta\), and let \(H=\operatorname{Stab}_G(B)\). Put
\(e=\dim B\), and suppose \(e\ge2\). The chosen block system
automatically satisfies (GSB.3), with an irreducible \(H\)-module
\(\tau=B\) and
\[
 \rho_{\mathbf C}=\operatorname{Ind}_H^G\tau,
 \qquad \mathbf Q(\chi_\tau)=K\mathbf Q(\chi_\theta).
 \tag{GSB.12}
\]
It therefore has the actual exact-field replacement and bounds
of (GSB.4)--(GSB.6). In a minimal source, \(H\) is normal,
\(G/H\) is the cyclic or Klein-four Galois group already described,
and \(\tau\) is faithful. In particular, if
\(\mathbf Q(\chi_\theta)\subseteq K\), then \(H=G\) and
\(\rho|_D\) is homogeneous.
\end{corollary}

\begin{proof}
Because \(D\) is normal, its canonical isotypic components are
permuted by \(G\). Irreducibility of \(\rho\) makes this action
transitive. They are also permuted by coefficient Galois action
fixing \(K\): that action commutes with the \(D\)-matrices and
permutes their irreducible character projectors. Thus this is the
promised automatic class of Galois-stable block systems.

The \(H\)-action on \(B\) is irreducible. Indeed a nonzero proper
\(H\)-invariant subspace would have direct translates in the
different \(D\)-isotypic components. Their sum would be a proper
nonzero \(G\)-subrepresentation. The translated direct sum of
the full \(B\)'s is \(\rho\), proving the induction in (GSB.12).

If \(\theta\) has multiplicity \(a\) in \(B\), then
\(\chi_\tau|_D=a\chi_\theta\), so \(\mathbf Q(\chi_\theta)\)
lies in \(L=\mathbf Q(\chi_\tau)\). We already know \(K\subseteq L\).
Conversely the central \(\theta\)-projector has coefficients in
\(\mathbf Q(\chi_\theta)\). Its image in \(\rho_K\) is defined
over \(K\mathbf Q(\chi_\theta)\), and is \(H\)-stable. The trace
of each \(H\)-matrix on that space lies in this compositum.
This proves the reverse field inclusion and hence (GSB.12).
Apply the theorem. If the compositum is just \(K\), then \(\ell=1\),
so minimality forces \([G:H]=1\). A single canonical component
is exactly homogeneity of the restriction.
\end{proof}

\paragraph{Scope retained.}
This reduces an automatic family of normal-subgroup systems; it
does not make every complex imprimitivity system Galois-stable.
If \(e=1\), the representation is monomial and the rational-base
argument for its inducing coefficient is not available here.
Only the separate monomial-source results, with their exact field
and nonsplit hypotheses, can remove that case.
A quasiprimitive representation may still be imprimitive; neither
the converse nor an exact-field reduction of every such source
is asserted. Proper \(L/K\) remains explicit throughout.
